% !TEX program = xelatex
\documentclass[12pt]{article}
\usepackage[margin=28mm]{geometry}
\usepackage{hyperref}
\usepackage{amsmath, amsthm, amssymb, amsfonts}
\usepackage{tikz-cd}
\usepackage{fontspec}
\usepackage{unicode-math}
\setmainfont{Latin Modern Roman}
\setmonofont{Latin Modern Mono}
\usepackage{listings}
\lstset{basicstyle=\ttfamily\footnotesize, columns=fullflexible, breaklines=true, showstringspaces=false}

\title{lean\_composition\_task — Lean→TeX Export (English)}
\author{TeX generated by ChatGPT-5 Pro\\Lean source generated by CODEX (OpenAI)}
\date{2025-09-10}

\begin{document}
\maketitle

\begin{abstract}
This document was automatically generated from a Lean source file. The author lines explicitly state: TeX generated by ChatGPT-5 Pro; Lean source generated by CODEX (OpenAI). The document includes a Provenance note, summary statistics, topic-appropriate diagrams (via tikz-cd), a table of first-line declarations, and the verbatim Lean source.
\end{abstract}

\section*{Provenance}
This \LaTeX{} file was generated from the Lean source file \texttt{lean\_composition\_task.lean}.\\
\textbf{TeX generation}: ChatGPT-5 Pro.\\
\textbf{Lean source generation}: CODEX (OpenAI).

\section*{Topic overview}
This note focuses on composition of linear maps U \xrightarrow{f} V \xrightarrow{g} W, highlighting linearity preservation, associativity, and identity laws. Proofs leverage standard lemmas such as \texttt{LinearMap.map_add} and function extensionality.

\section*{At-a-glance}
Total lines: 69\\
Defs: 0\\
Lemmas: 0\\
Theorems: 3\\
Structures: 0\\
Inductives: 0

\section*{Sample diagrams}

\[
\begin{tikzcd}
A \arrow[r, "f"] \arrow[d, "g"'] & B \arrow[d, "h"] \\
C \arrow[r, "k"'] & D
\end{tikzcd}
\]



\[
\begin{tikzcd}
U \arrow[r, "f"] \arrow[rr, bend left=15, "{g \circ f}"] & V \arrow[r, "g"] & W
\end{tikzcd}
\]


\section*{Declarations (first-line signatures)} 
\begin{center}
\begin{tabular}{lll}
\textbf{kind} & \textbf{name} & \textbf{signature (first line)} \\ \hline
theorem & \texttt{comp\_preserves\_linearity} & \texttt{{R U V W : Type*} [Semiring R] [AddCommMonoid U] [AddCommMonoid V] [AddCommMonoid W] [Module R U] [Module R V] [Module R W] (g : V →ₗ[R] W) (f : U →ₗ[R] V) (u₁ u₂ : U) : (g.comp f) (u₁ + u₂) = (g.comp f) u₁ + (g.comp f) u₂ := by} \\
theorem & \texttt{comp\_assoc} & \texttt{{R U V W X : Type*} [Semiring R] [AddCommMonoid U] [AddCommMonoid V] [AddCommMonoid W] [AddCommMonoid X] [Module R U] [Module R V] [Module R W] [Module R X] (h : W →ₗ[R] X) (g : V →ₗ[R] W) (f : U →ₗ[R] V) : (h.comp g).comp f = h.comp (g.comp f) := by} \\
theorem & \texttt{comp\_id\_left} & \texttt{{R V W : Type*} [Semiring R] [AddCommMonoid V] [AddCommMonoid W] [Module R V] [Module R W] (f : V →ₗ[R] W) : LinearMap.id.comp f = f := by} \\
\end{tabular}
\end{center}

\section*{Lean source (verbatim)}
\begin{lstlisting}
import Mathlib.Algebra.Module.LinearMap.Defs
import Mathlib.Algebra.Module.LinearMap.Basic
import Mathlib.Algebra.Module.Prod

/-!
  Bourbaki の精神（順序対と射影）：
  構造は順序対とその射影 π₁, π₂ により特徴付けられる。線形写像の合成も、
  値を順序対に入れて射影で読むことで成分ごとに確認できる。
-/

noncomputable section

namespace MyProjects.LinearAlgebra.A

-- 線形写像の合成の線形性
theorem comp_preserves_linearity {R U V W : Type*}
  [Semiring R]
  [AddCommMonoid U] [AddCommMonoid V] [AddCommMonoid W]
  [Module R U] [Module R V] [Module R W]
  (g : V →ₗ[R] W) (f : U →ₗ[R] V) (u₁ u₂ : U) :
  (g.comp f) (u₁ + u₂) = (g.comp f) u₁ + (g.comp f) u₂ := by
  simpa using (LinearMap.map_add (g.comp f) u₁ u₂)

-- 線形写像の合成の結合法則
theorem comp_assoc {R U V W X : Type*}
  [Semiring R]
  [AddCommMonoid U] [AddCommMonoid V] [AddCommMonoid W] [AddCommMonoid X]
  [Module R U] [Module R V] [Module R W] [Module R X]
  (h : W →ₗ[R] X) (g : V →ₗ[R] W) (f : U →ₗ[R] V) :
  (h.comp g).comp f = h.comp (g.comp f) := by
  ext u; rfl

-- 発展：恒等写像との合成
theorem comp_id_left {R V W : Type*}
  [Semiring R] [AddCommMonoid V] [AddCommMonoid W]
  [Module R V] [Module R W]
  (f : V →ₗ[R] W) :
  LinearMap.id.comp f = f := by
  ext v; rfl

/-! ## Bourbaki スタイルの射影による成分チェック（例） -/

section Examples
  variable {R U V W : Type*}
  variable [Semiring R]
  variable [AddCommMonoid U] [AddCommMonoid V] [AddCommMonoid W]
  variable [Module R U] [Module R V] [Module R W]

  -- 合成の加法性を値レベル（関数合成として）で確認
  example (g : V →ₗ[R] W) (f : U →ₗ[R] V) (u₁ u₂ : U) :
      g (f (u₁ + u₂)) = g (f u₁) + g (f u₂) := by
    simpa using (LinearMap.map_add (g.comp f) u₁ u₂)

  -- 順序対値の線形写像に対して、射影で成分ごとに `map_add` を読む
  example (f : U →ₗ[R] V × W) (u₁ u₂ : U) :
      (f (u₁ + u₂)).1 = (f u₁).1 + (f u₂).1 := by
    simpa using congrArg Prod.fst (LinearMap.map_add f u₁ u₂)

  example (f : U →ₗ[R] V × W) (u₁ u₂ : U) :
      (f (u₁ + u₂)).2 = (f u₁).2 + (f u₂).2 := by
    simpa using congrArg Prod.snd (LinearMap.map_add f u₁ u₂)

  -- 結合律は成分に適用しても同じ（`ext` での等式確認）
  example (h : W →ₗ[R] W) (g : V →ₗ[R] W) (f : U →ₗ[R] V) :
      ((h.comp g).comp f) = h.comp (g.comp f) := by
    ext u; rfl
end Examples

end MyProjects.LinearAlgebra.A

\end{lstlisting}

\end{document}
